% Handoff notes and table for manuscript revision.
% LaTeX fragment; requires amsmath and booktabs in the parent document.
% These are NOT results of the live 250-case-per-task RetinAgent experiment.

\subsection*{Scope of the available escalation analysis}

The available curves describe an earlier deterministic, reliability-weighted
multimodel fusion experiment on Harvard FairVision30K. They do not describe
the current live OphthalmicAgent/RetinAgent pipeline. The source contains
3,000 held-out test cases for each task: 1,466 positive and 1,534 negative
glaucoma cases; 1,069 positive and 1,931 negative AMD cases; and 261 positive
and 2,739 negative DR cases. These cohorts retain their observed task-specific
prevalence and are distinct from the 250-case diagnostic evaluation subsets.

The historical source lists nine model--modality configurations: RETFound
OCT, MIRAGE SLO, FLAIR SLO, RET-CLIP SLO, VisionFM OCT and SLO, RetiZero SLO,
and UrFound OCT and SLO. Its recorded reliability formula is
\begin{equation}
s = 0.45\,\mathrm{AUROC} + 0.35\,\mathrm{F1}
    - 0.10\,\mathrm{ECE} - 0.07\,\mathrm{FNR} - 0.03\,\mathrm{FPR}.
\end{equation}
This is historical configuration information, not a replacement for the
current manuscript's reliability formula. The present plotting audit does
not establish the provenance or split purity of every upstream reliability
estimate. It checks and summarizes the saved predictions and risk scores.

\subsection*{How the curves were computed}

For case $i$, let $p_i$ be the saved fused probability, $v_i$ the fraction of
models voting positive, and $q_i$ the saved weighted reliability. The recorded
deferral score is
\begin{equation}
d_i = \max\left\{1-2|p_i-0.5|,\;
                  \min(v_i,1-v_i),\;
                  \max(0,0.35-q_i)\right\}.
\end{equation}
Larger scores indicate stronger grounds for deferral in this historical
scoring rule; they are not calibrated probabilities of error. The saved
classification label is $\hat y_i=\mathbf{1}[p_i\geq0.5]$ and is held fixed
throughout the analysis. At threshold $\tau$, the accepted set is
$A_\tau=\{i:d_i\leq\tau\}$. Coverage is $|A_\tau|/N$ and the percentage
escalated is $100(1-|A_\tau|/N)$.

The figure sweeps every distinct saved risk score, retaining equal scores
together. Lines connect feasible operating points without smoothing; points
between them do not represent additional measured thresholds. The table shows
fixed illustrative score thresholds, not clinically selected operating
points. This analysis does not fit new predictions, optimize a threshold,
or make API calls. It does not reapply the original joint policy involving
conformal prediction, close calls and disagreement flags.

On each accepted subset, macro-F1 is recomputed as
\begin{equation}
F_{1,\mathrm{macro}}=
\frac{1}{2}\left(
\frac{2\mathrm{TP}}{2\mathrm{TP}+\mathrm{FP}+\mathrm{FN}}+
\frac{2\mathrm{TN}}{2\mathrm{TN}+\mathrm{FP}+\mathrm{FN}}
\right).
\end{equation}
Both labels are fixed as $\{0,1\}$; a zero class-wise denominator contributes
zero, while an empty accepted subset has undefined F1. The original saved
risk--coverage CSV reported positive-class F1, so those values were not
relabelled as macro-F1. Single-class subsets remain in the exported source
data but are not plotted. Dashed curve segments mark fewer than 20 accepted
positive or 20 accepted negative cases; this is a display warning, not a
threshold-selection rule.

\begin{table*}[t]
\centering
\caption{Exploratory accepted-case performance of historical deterministic
reliability-weighted fusion on Harvard FairVision30K, not the live RetinAgent
evaluation. Each task includes 3,000 held-out cases. Cases with saved risk
score $d_i\leq\tau$ are accepted; all other cases are designated for review.
Classification labels remain fixed. Thresholds are illustrative and were
not optimized in this analysis. F1 is macro-averaged over accepted cases;
positive/negative counts refer to their reference labels.
$\dagger$ marks fewer than 20 accepted cases in either class.
At $\tau=0$, all cases are escalated and accepted-case F1 is undefined.}
\label{tab:historical_fusion_escalation}
\small
\begin{tabular}{@{}lrrrrrc@{}}
\toprule
Task & $\tau$ & Accepted $n$ & Coverage (\%) & Escalated (\%) &
Accepted $+/-$ & Macro-F1 \\
\midrule
Glaucoma & 1.0 & 3,000 & 100.00 &  0.00 & 1,466/1,534 & 0.7568 \\
         & 0.8 & 2,476 &  82.53 & 17.47 & 1,236/1,240 & 0.8063 \\
         & 0.6 & 1,732 &  57.73 & 42.27 &   915/817   & 0.8643 \\
         & 0.4 &   797 &  26.57 & 73.43 &   486/311   & 0.9297 \\
         & 0.2 &    85 &   2.83 & 97.17 &    83/2     & $1.0000^{\dagger}$ \\
\midrule
AMD      & 1.0 & 3,000 & 100.00 &  0.00 & 1,069/1,931 & 0.8022 \\
         & 0.8 & 2,495 &  83.17 & 16.83 &   926/1,569 & 0.8551 \\
         & 0.6 & 1,833 &  61.10 & 38.90 &   792/1,041 & 0.9013 \\
         & 0.4 & 1,046 &  34.87 & 65.13 &   602/444   & 0.9505 \\
         & 0.2 &   355 &  11.83 & 88.17 &   304/51    & 0.9639 \\
\midrule
DR       & 1.0 & 3,000 & 100.00 &  0.00 &   261/2,739 & 0.7398 \\
         & 0.8 & 2,695 &  89.83 & 10.17 &   191/2,504 & 0.7965 \\
         & 0.6 & 2,440 &  81.33 & 18.67 &   121/2,319 & 0.8191 \\
         & 0.4 & 1,954 &  65.13 & 34.87 &    47/1,907 & 0.8418 \\
         & 0.2 & 1,170 &  39.00 & 61.00 &    17/1,153 & $0.8689^{\dagger}$ \\
\bottomrule
\end{tabular}
\end{table*}

\subsection*{Interpretation and manuscript wording}

In this exploratory analysis, at an illustrative score threshold of 0.8,
accepted-case macro-F1 was 0.8063, 0.8551 and 0.7965 for glaucoma, AMD and DR,
respectively, compared with full-coverage values of 0.7568, 0.8022 and 0.7398.
The corresponding proportions designated for clinician review were 17.47\%,
16.83\% and 10.17\%. At a threshold of 0.4, accepted-case macro-F1 was higher
(0.9297, 0.9505 and 0.8418), but review rates rose to 73.43\%, 65.13\% and
34.87\%, respectively. These values describe a trade-off between retained
coverage and performance within the retained subset, not improved predictions
for the entire cohort.

The full curves need not improve monotonically as coverage falls. Their
extreme tails have limited support: glaucoma reaches macro-F1 of 1.0000 at
$\tau=0.2$ while retaining only 85 cases, of which just two are negative.
For DR, the same threshold retains only 17 positive cases among 1,170 accepted
cases. F1 depends on class composition, which changes with the threshold.
Higher accepted-case macro-F1 therefore cannot be attributed solely to better
uncertainty ranking. No uncertainty intervals were computed for these curves,
and the separate 250-case glaucoma confidence intervals do not apply to them.
There are no clinician outcomes for the escalated cases, and designating cases
for review does not establish that they would be correctly diagnosed or that
overall clinical workload would decrease.

If included, this result belongs under a clearly identified exploratory
deterministic-fusion analysis, preferably supplementary. It should not fill
the main-text claim that the live RetinAgent system achieves a particular
accepted-case macro-F1 or escalation rate. It also does not establish that
agentic reasoning, the current reliability formula, or subgroup conditioning
caused an improvement. No best threshold or deployable operating point is
claimed from this post-hoc test-cohort analysis.

\subsection*{What remains necessary for the live RetinAgent figure}

For each locked task cohort, obtain one record per case containing the case
identifier, reference label, final agent prediction and a saved, well-defined
deferral score linked to that same prediction and run. Reference labels are
used for evaluation only, not for constructing the score. Failed or missing
cases must be resolved or handled under a declared rule, not silently omitted.
A binary escalation flag supports one operating point, not a continuous
threshold sweep. Unlinked counterfactual confidence reports should not be
substituted for final-output confidence. If no continuous score was saved,
report the existing operating point or evaluate an explicitly defined new
deferral policy; do not describe such a policy as the original agent's rule.
Any new operating threshold should be selected on development data and then
evaluated on an untouched cohort.

% Sources within this directory: threshold_checkpoints.csv,
% threshold_curve.csv and provenance.json. Figure: historical_fusion_escalation_macro_f1.pdf.
% Original case predictions:
% equi-agent/outputs/fairvision_reliability_selective_arbitration/selective_arbitration_predictions.csv
